\subsection{Grundverhalten eines npn-Bipolartransistors}
% Alt: Praktikumsaufgabe 8
\Aufgabe{
    \label{Aufg35}
    Analysieren Sie das Grundverhalten eines Bipolartransistors in einer Basisschaltung.
    Der npn-Transistor weist einen Stromverstärkungsfaktor von $\beta = 100$ auf,
    es wird ein Basisstrom von $I_\mathrm{B} = 20\,\mathrm{\mu A}$ angenommen.

    \begin{figure}[H]
        \centering
        \input{Tikz/src/FigPraktikumNPN}\alttext{Das Bild zeigt ein elektrisches Schaltbild mit einem Transistor. Links befindet sich eine Spannungsquelle \(U_1\), die mit zwei Widerständen \(R_1\) und \(R_2\) in Reihe geschaltet ist. Der Verbindungspunkt zwischen \(R_1\) und \(R_2\) führt über einen Widerstand \(R_3\) zum Basisanschluss des Transistors. Vom Kollektor des Transistors geht ein Widerstand \(R_4\) ab, der zu einer zweiten Spannungsquelle \(U_2\) führt. Der Emitter des Transistors ist mit Masse verbunden. Am Widerstand \(R_3\) ist der Basisstrom \(I_\mathrm{B}\) mit einem roten Pfeil markiert. Die Stromrichtungen der Spannungsquellen \(U_1\) und \(U_2\) sind jeweils durch blaue Pfeile nach unten gekennzeichnet.}
        \label{fig:FigPraktikumNPN}
    \end{figure}


    \begin{itemize}
        \item[a)] Berechnen Sie den Kollektorstrom $I_\mathrm{C}$.
        \item[b)] Bestimmen Sie die Kollektor-Emitter-Spannung $U_\mathrm{CE}$, wenn der Kollektorwiderstand $R_\mathrm{4} = 1\,\mathrm{k\Omega}$ und die Versorgungsspannung $U_\mathrm{2} = 12\,\mathrm{V}$ beträgt.
    \end{itemize}
}


\Loesung{
    \begin{itemize}
        \item[a)] Berechnung des Kollektorstroms $I_\mathrm{C}$
              Der Kollektorstrom eines Bipolartransistors ergibt sich aus der Beziehung:
              $$ I_\mathrm{C} = \beta \cdot I_\mathrm{B} $$
              mit den gegebenen Werten:
              \begin{equation*}
                  I_\mathrm{C} = 100 \cdot 20\,\mathrm{\mu A} = 2\,\mathrm{mA}
              \end{equation*}


        \item[b)] Bestimmung der Kollektor-Emitter-Spannung $U_\mathrm{CE}$
              Die Kollektor-Emitter-Spannung ergibt sich aus der Gleichung:
              $$ U_\mathrm{CE} = U_\mathrm{CC} - I_\mathrm{C} \cdot R_\mathrm{4} $$
              Einsetzen der Werte:
              \begin{equation*}
                  U_\mathrm{CE} = 12\,\mathrm{V} - (2\,\mathrm{mA} \cdot 1\,\mathrm{k\Omega}) = 12\,\mathrm{V} - 2\,\mathrm{V}  = 10\,\mathrm{V}
              \end{equation*}

    \end{itemize}
    Siehe: Abschnitt \ref{fig:UebersichtBipolartransistoren}
}